\section{Introduction}\label{met:introduction}
Countably branching diamond graphs test whether a Banach space can
accommodate many separated approximate midpoints at every scale. Each
edge is replaced by countably many two-edge paths with common endpoints.
A target space that forces these midpoint families away from the sphere
must lose some edge stretch at every replacement. The question is what
quantity decreases, and how the resulting distortion bound depends on
the geometry used to prove that decrease.

We use the normalized diamond $\mathcal D_k$: the original edge has
length one, and the two edges in each replacement path have half the
length of their parent. Thus $\mathcal D_k$ has finite depth and
countably many vertices; write $d_k$ for its shortest-path metric. A map $f$ has distortion at most $D$ if some
$s>0$ satisfies
\[
 s d_k(u,v)\le\|f(u)-f(v)\|\le Ds d_k(u,v)
 \qquad(u,v\in\mathcal D_k).
\]
The scale $s$ may depend on the map and on $k$. Its removal by
multiplying $f$ by a scalar does not change $D$.

Our main example applies James's disjoint-segment norm
\cite[p.~739]{James1974} to a forest of finite heights. Let
$\mathcal T=\coprod_{h\ge1}\mathbb N^{\le h}$, with the prefix order
within each rooted tree. A segment is the set of consecutive vertices
between two comparable nodes. Complete the finitely supported real
arrays on $\mathcal T$ in the norm
\[
 \rho(a)=\sup\left\{
  \left(\sum_{j=1}^m\left|\sum_{s\in S_j}a_s\right|^2\right)^{1/2}:
  S_1,\ldots,S_m\text{ are pairwise disjoint segments}\right\},
\]
and call the resulting space $J$. We consider $X=J^*$ with its given
dual norm. On the infinite tree $T_\infty=\mathbb N^{<\omega}$,
including its root, let $J_\infty$ be the completion of the finitely
supported real arrays for the same disjoint-segment formula. Define
\[
 B_\infty=\overline{\operatorname{span}}\{e_t^*:t\in T_\infty\}
 \subset J_\infty^*,
\]
where $e_t^*$ is coordinate evaluation and the closure uses the dual
norm. The companion on midpoint lenses proves the midpoint estimates
and the finite-coordinate approximation facts for these two spaces
\cite{SegmentLenses}. Section~\ref{dir:forest-definition} also gives
the forest foundations needed for our direct proofs.

James's original construction uses the infinite binary tree.
Girardi proved asymptotic uniform convexity for the coordinate predual
and the dual of that binary-tree space
\cite[Theorems~3 and~5]{Girardi2001}. The countably branching spaces
defined here use their own segment estimates.

\begin{theorem}\label{met:intro-main}\label{met:joint-distortion}
For every integer $k\ge0$, every distortion-$D$ embedding of
$\mathcal D_k$ into either $X$ or $B_\infty$ satisfies
\[
 D^2\ge1+k/4.
\]
\end{theorem}

In both cases the input is the following tail inequality. A finite set
$H$ is ancestral if it contains every ancestor of each of its vertices;
$P_H$ retains its coordinates and $Q_H=I-P_H$. If $P_Hp=p$, then
\[
 \|p\pm y\|\le R
 \quad\Longrightarrow\quad
 \|Q_Hy\|\le2\sqrt{R^2-\|p\|^2}.
\]
It holds for arbitrary displacements $y$. In a symmetric lens
$\{z:\|x\pm z\|\le R\}$ containing infinitely many separated deviations,
two projected deviations can be made arbitrarily close, and their
half-difference remains in the same lens. This turns the tail inequality
into a loss in squared edge stretch. Section~\ref{met:joint-section}
gives the complete deduction. The coefficient $1/4$ results from retaining the full quadratic budget
in this calculation.

\subsection{Midpoint methods and earlier work}
Diamond graphs connect metric embeddings with convexity of Banach-space
norms. For binary diamonds, Lee and Naor obtain
quantitative lower bounds in $L_p$, $1<p\le2$, by summing a convexity
inequality over the replacement levels \cite[Lemma~2.2]{LeeNaor2004}.
Johnson and Schechtman characterize superreflexivity by the failure of
uniformly bounded binary-diamond embeddings; their lower-bound proof
iterates a modulus-of-convexity estimate
\cite[Section~4]{JohnsonSchechtman2009}. The countably branching graphs
studied here make a further selection possible: infinitely many midpoint
images can be compared after their projections onto a fixed
finite-dimensional space have become arbitrarily close.

Dilworth, Kutzarova, Randrianarivony, Revalski, and Zhivkov introduced
asymptotic midpoint uniform convexity and characterized it by uniform
smallness of the noncompactness of symmetric approximate-midpoint sets
\cite[Theorem~2.1 and Definition~2.2]{DKRRZ2016}.
In finite-codimensional terms, for each $t>0$ there is $\delta>0$
such that every unit center $x$ admits a finite-codimensional subspace
on whose unit directions $y$ one has
$\max\{\|x+ty\|,\|x-ty\|\}\ge1+\delta$. Replacing the maximum
by $\|x+ty\|$ gives asymptotic uniform convexity.
Baudier, Causey, Dilworth, Kutzarova, Randrianarivony, Schlumprecht, and Zhang turn this property into a countably branching
diamond obstruction \cite[Section~4]{BaudierEtAl2017}. Their
Theorem~4.1 includes a quantitative conclusion: a midpoint modulus bounded
below by a positive multiple of $t^p$ for $0<t<1$, with
$1<p<\infty$, gives distortion
bounded below by a positive multiple of $k^{1/p}$. Their proof combines control of midpoint
sets with an iterated improvement of endpoint distances. The slope-loss
arguments here use that mechanism and retain the constants furnished by
the stated vector estimates. Malthaner gives another proof of the
power-type lower bound through coloring, a path lemma, and separated
approximate-midpoint configurations
\cite[Theorem~3.5 of arXiv v2]{Malthaner2025}. He also proves a power-type obstruction for metric targets
satisfying a diamond inequality
\cite[Equation~(2) and Theorem~3.7 of arXiv v2]{Malthaner2025}.
The squared-lens inequality used here is its $p=2$ Banach-space form;
our coefficient comes from the stated norm-specific tail estimate.

Converse embedding results use additional geometry. For reflexive spaces
with unconditional asymptotic structure, Baudier et al. characterize
existence of an equivalent asymptotically uniformly convex norm by the
failure of uniformly bounded diamond embeddings
\cite[Corollary~5.3]{BaudierEtAl2017}.
Swift extends the characterization to countably branching bundle graphs
with these hypotheses \cite[Corollary~3.4 and Theorem~6.9]{Swift2018}.
Perreau gives a dual characterization for separable $Y$ when $Y^*$ has
weak-star unconditional asymptotic structure, and a characterization for
reflexive spaces with unconditional asymptotic structure
\cite[Theorems~4.8 and~5.3]{Perreau2021}. These assumptions remain part
of the corresponding converses.

Normalized diamond slopes also occur in the sprawling-tree construction
of Basset, Lancien, and Proch\'azka
\cite[Definition~3.7 and Proposition~6.4]{BassetLancienProchazka2026}.
Our direct proofs specialize that construction to finite height and then
supply coordinate selections, primal corrections, and predictions of
atomic segment coefficients for the norms considered here.

We distinguish numerical conventions for the midpoint modulus.
Dilworth et al.\ \cite[Definition~2.2]{DKRRZ2016} and
Baudier et al.\ \cite[Section~4]{BaudierEtAl2017} use the maximum of the
two signed norms. The Baudier--Lancien book manuscript uses their
average \cite[Definition~5.1.15 of arXiv v1]{BaudierLancien2025}.
An average lower bound implies the corresponding maximum lower bound;
each application below specifies which estimate it uses.

\subsection{Direct proofs and the organization of the paper}
The forest also admits direct graph arguments that construct their own
energy budgets. These proofs use the segment norm itself, without invoking
the companion's midpoint inequality. The common first step is to select
one replacement path at each retained edge. Infinite branching permits
coordinate choices at this step; subsequent descent uses only the two
half-edges of each selected path.

Sections~\ref{dir:forest-definition} and~\ref{dir:accumulated} develop
the forest foundations and two primal-test constructions. A test is a
finitely supported vector in $J$, evaluated by the current slope in $X$.
One construction adds a correction at the first vertices outside a finite
head, so that a gain in evaluation costs only a bounded increase in the
squared test norm. The other balances prefix sums at each dyadic size and
bounds the final accumulated test by charging segment sums to their
largest encountered size. Both compare growing evaluations with the
norm of the test.

Section~\ref{dir:atomic:section} uses a different budget. Fair choices
between the two half-edges make the selected slopes a finite martingale.
After projection to a common finite head, the terminal slopes are
represented exactly by random sums of disjoint-segment functionals whose
squared coefficients sum to at most one. As more of a
segment's location is revealed, conditional predictions of its coefficient
have bounded total squared change. Counting segment endpoints, or only
segment starts, finds an interval of levels where this prediction energy
cannot sustain all the large new coordinate layers selected from the graph.
The coordinate selections and test estimates here are specific to our norm.

Section~\ref{dir:layers} then treats the linear path-cost norm. It
states each imported energy estimate alongside the graph construction
that meets its support and timing hypotheses. Some constructions keep the
exact graph slopes and allow small tails outside a head; others replace
the slopes by finitely supported approximations while preserving their
exact averaging identity.

Section~\ref{met:transfer-section} collects the remaining midpoint
transfers after these direct proofs.
They distinguish squared or higher-power losses, a fixed contraction from
a separated lens, and growth of an actual child slope.
Sections~\ref{met:d-section} and~\ref{met:f:applications} introduce the
other path-cost, segment-start and recursive-array norms and apply these
transfers. Each application states its intrinsic input from the complete
companions on bounded tree potentials or recursive potentials. Linear
and quadratic path costs, free and fixed joining roots, and different
child-aggregation exponents give different norms; the stated constants
belong to those specified norms and inputs.
Appendix~\ref{met:forest-variants} records the alternative forest
midpoint arguments and their individual quantitative or qualitative losses.
